\documentclass[11pt,a4paper]{article}
\usepackage[utf8]{inputenc}
\usepackage[T1]{fontenc}
\usepackage[spanish]{babel}
\usepackage{hyperref}
\usepackage{geometry}
\geometry{margin=2.5cm}
\usepackage{parskip}

\title{107 el md es para leer y el jsonl es para probar}
\author{Alejandro Valenzuela Rivera\\\small Universidad Autónoma de Nuevo León}
\date{\today}
\begin{document}
\maketitle

\begin{quote}
\textbf{Sesgo declarado.} Este trabajo se sostiene desde el sesgo declarado de su autor, y no finge neutralidad. Lo que afirma se puede tumbar, y ese es el punto: el falsificador está escrito, no escondido.
\end{quote}

\section*{Resumen}
Un chat con tool calls no se pierde en md si se separa: el md lleva la conversacion y las llamadas legibles, y el jsonl lleva el evento crudo con su uso. Lo que no se puede es meter el crudo en el md ni el md en el crudo.

\section{Por qué se sostiene}
Y la cuenta lo dice: esta sesion tiene 1004 tool call y 1024 tool result contra 305 mensajes del autor. O sea, la parte mas interesante es la que no es conversacion, y si se aplana a prosa se pierde justo lo que sirve de prueba. Por eso el md es la lectura y el jsonl es la prueba.

\section{De dónde viene}
La sesion medida hoy.

\section{El caso}
Ninguna: es medicion.

\section{Método}
Un .md que se lee y un .jsonl que se corre, y cada uno hace lo suyo.

\section{Límites, y qué lo tumba}
Un .md usado para probar, o un .jsonl del que se saca una conclusion a ojo.

\section{Lo que queda abierto}
\textquestiondown{}Si el md pierde las tool calls?

\section*{Cómo se leyó esta tesis}
Esto se escribió en el repositorio man-made-empathy, del que este candado es una pieza: candados/04 el lenguaje/107 el md es para leer y el jsonl es para probar.md. La máquina que lo compuso lo hizo desde este mismo archivo, y no agregó nada que no estuviera escrito. Si algo de esto es falso, se dice en el mismo sitio, y se corrige con fecha.

\end{document}
